% [原文 p.257]
\chapter{半稳定极小模型}

本章解说称为半稳定 MMP 的、在特殊情形下的 MMP。与其说处理三维簇本身，不如把它视为曲线上的曲面族。半稳定性是一个技术性假设：这些曲面的退化不能太复杂。在此假设下，我们证明三维翻转的存在性，特别是 MMP 能够运行。

[Mor88] 关于三维翻转存在性的证明，以及 [Sho92] 更为一般的处理，都很长且复杂。与此相对，半稳定翻转虽然特殊，却仍具备一般翻转的许多有趣特征。此外，半稳定 MMP 还有若干有趣的应用。作为三维翻转分类 [KM92] 的结果，我们知道几乎所有翻转都是半稳定的，但直接证明这一点非常困难。

§7.1 解说半稳定 MMP 的一般设定。

§7.2 对 [KKMSD73] 的半稳定约化定理在三维情形给出证明。这里不使用环面嵌入理论，而证明第 6 章所建立的 Q-分解终端化就是半稳定奇点消解。

§7.3 与 §7.4 讨论半稳定翻转。首先在 §7.3 证明所谓特殊半稳定翻转的存在性，然后在主要部分 §7.4 把任意半稳定翻转的存在性归约到它。那里从半稳定翻转收缩出发，作辅助性操作。于是，经过只由特殊半稳定翻转组成的
% [原文 p.258]
另一个半稳定 MMP，就得到原目标的半稳定翻转。这一手法最初由 [Sho92] 在稍有不同的情形下使用。这里所用的论证基于 Corti 的想法。

§7.5 处理半稳定 MMP 的 3 个应用。无论哪种情形，证明的核心都由半稳定 MMP 承担，我们详细解说这一部分。证明的其他部分基于模空间、奇点理论等已经确立的方法。这些部分就留给各自的相关文献。

最后，§7.6 是关于最近发展的概说。

\section{半稳定极小模型纲领}
\takeaway{本节先以 lc 射与 dlt 射给出半稳定三维 MMP 的设定，再叙述其两大支柱：半稳定翻转的存在（§7.4）与半稳定翻转的终止（本节末）。}

MMP 依所适用的情形而略有不同形态。本节定义半稳定三维 MMP，并给出其能够运行的全部证明：半稳定翻转的存在（推论 7.42）在 §7.4 证明，半稳定翻转的终止（定理 7.7）在本节末尾证明。

我们从 lc 射与 dlt 射的定义开始，用它们可以高效地定义半稳定 MMP。

\begin{definition}{7.1}
设 $X$ 是正规簇，$B=\sum b_iB_i$ 是有效 $\Q$-除子，且 $f\colon X\to C$ 是到非奇异曲线 $C$ 的非常数态射。若对每个闭点 $c\in C$，$(X,B+f^{-1}(c))$ 是 lc（或 dlt）的，则称态射 $f\colon(X,B)\to C$ 为对数标准（log canonical）的（或除子对数终端（divisorial log terminal）的，简称 lc（或 dlt））。
\end{definition}

% [Tutor Note: lc 射要求每个闭点处的对 $(X,B+f^{-1}(c))$ 都 lc，于是由定义 2.34 其分支系数不超过 1，纤维自动约化、$B$ 自动水平；这正是把逐纤维的 lc 性打包成一个射的条件的好处。]

\begin{lemmap}{7.2}
假设射 $f\colon(X,B)\to C$ 是 lc 的。则以下成立。
\begin{enumerate}[label=(\arabic*),leftmargin=2.4em,itemsep=1pt,topsep=2pt,parsep=0pt]
\item $f$ 的任意纤维都是约化的。
\item $B$ 是水平的，即 $B$ 的任意不可约分支都不含于 $f$ 的纤维中。
\item $(X,B)$ 是 lc 的。
\item 若 $E$ 是 $X$ 上方的任意例外除子，且 $\operatorname{center}_XE$ 含于某纤维中，则 $a(E,X,B)\ge0$。
% [原文 p.259]
\item $(X,B)$ 是标准的（或 klt）当且仅当 $(X_{gen},B|_{X_{gen}})$ 是 lc（或 klt）的，其中 $X_{gen}$ 是 $f$ 的一般纤维。
\end{enumerate}
\end{lemmap}

\begin{proof}
取 $c\in C$。由定义 2.34，$B+f^{-1}(c)$ 的任意不可约分支的系数都不超过 1。这意味着 $f^{-1}(c)$ 是约化的，且 $B$ 与 $f^{-1}(c)$ 没有公共分支。由推论 2.35，$(X,B)$ 是 lc 的。

假设 $\operatorname{center}_XE\subset f^{-1}(c)$。此时
\[
a(E,X,B)\ge a(E,X,B+f^{-1}(c))+1\ge0
\]
成立。设 $E$ 是 $X$ 上方的例外除子，且 $\operatorname{center}_XE$ 不含于任何纤维中，用 $E_{gen}$ 表示 $E\to C$ 的一般纤维。$E_{gen}$ 的任意不可约分支 $E'_{gen}$ 是 $X_{gen}$ 的素除子。我们来证明 $a(E,X,B)=a(E'_{gen},X_{gen},B|_{X_{gen}})$（这样便得到 (5)）。事实上，若 $E$ 是双有理态射 $h\colon Z\to X$ 上的除子，则把
\[
K_Z=h^*(K_X+B)+a(E,X,B)E+\cdots
\]
限制到 $Z_{gen}$，得到 $K_{Z_{gen}}=h_{gen}^*(K_{X_{gen}}+B_{gen})+a(E,X,B)E'_{gen}+\cdots$，即所证之主张成立。
\end{proof}

\begin{remark}{7.3}
尽管引理 7.2(4) 给人以某种印象，$(X,B)$ 一般并不是标准的。例如，若 $S$ 是任意 lc 曲面，则 $X=C\times S\to C$ 是 lc 的，而 $X$ 标准当且仅当 $S$ 标准。尽管如此，在许多方面 $(X,B)$ 表现得像标准对。

如果在某处需要用一个对标准对成立、但对 lc 对不成立的定理，值得把证明过一遍。稍作修改，它往往对 lc 射也能奏效。
\end{remark}

% [Tutor Note: (5) 把小维数现象一般化：标准性与 klt 性只取决于一般纤维，而 lc 性还允许沿纤维出现系数恰为 1 的分支。引理 7.4 就是曲面情形下这一差别的具体形态。]

\begin{lemmap}{7.4}
假设 $f\colon(X,B)\to C$ 是 lc 射，$c\in C$，$x\in f^{-1}(c)$ 是闭点。
\begin{enumerate}[label=(\arabic*),leftmargin=2.4em,itemsep=1pt,topsep=2pt,parsep=0pt]
\item 若 $\dim X=2$，则下列之一成立。
\begin{enumerate}[label=(\alph*),leftmargin=2.2em,itemsep=1pt,topsep=1pt,parsep=0pt]
\item $f^{-1}(c)$ 在 $x$ 处有通常结点，且 $x\notin\operatorname{Supp}B$。
\item $f^{-1}(c)$ 在 $x$ 处非奇异，且 $\operatorname{mult}_xB\le1$。
\end{enumerate}
\item 若 $f^{-1}(c)$ 在 $x$ 处正规，则 $(f^{-1}(c),B|_{f^{-1}(c)})$ 在 $x$ 处是 lc 的。
\end{enumerate}
\end{lemmap}

\begin{proof}
若 $\dim X=2$，则由引理 7.2(4) 知 $(X,B)$ 是标准的，再由定理 4.5 得到 (1) 的其余部分。(2) 由命题 5.46 得到。
\end{proof}

% [原文 p.260]
\begin{remark}{7.5}
许多应用中，理解 $f^{-1}(c)$ 的非正规奇点也很重要。正是这个问题引出了 lc 对概念的非正规版本。所谓半对数标准（简称 slc）对的理论与 lc 对的情形相差不大，但仍须解决一些基本问题。详见 [KSB88, K+92]。
\end{remark}

\begin{lemmap}{7.6}
假设射 $f\colon(X,B)\to C$ 是 lc 的。设 $g\colon C'\to C$ 是从非奇异曲线 $C'$ 出发的非常数态射，$X':=X\times_CC'$，而 $h\colon X'\to X$ 与 $f'\colon X'\to C'$ 为投影。记 $B':=h^*B$（参看命题 5.20）。

此时，射 $f'\colon(X',B')\to C'$ 也是 lc 的。
\end{lemmap}

\begin{proof}
取包含 $g$ 的全部分歧点的有限集 $R\subset C$，并记 $R':=\operatorname{red}(g^{-1}(R))$。我们证明 $X'$ 是正规的，且下式成立：
\[
K_{X'}+f'^{-1}(R')+B'=h^*(K_X+f^{-1}(R)+B).
\]
由于 $C'\to C$ 是平坦的，$X'\to X$ 也平坦。由于 $X$ 是 $S_2$ 的，这意味着 $X'$ 也是 $S_2$ 的（命题 5.4）。因此，$X'$ 只要在余维 1 的点处正规，就是正规的。

由于 $f$ 的所有纤维都是约化的，$f$ 在余维 2 的子集 $W\subset X$ 之外是光滑的。于是 $f'$ 也在余维 2 的子集 $W':=h^{-1}(W)$ 之外光滑。由此，$X'$ 在所有余维 1 的点处正规，且 $f'$ 的所有纤维都是约化的。

这也说明 $K_{X'/C'}=h^*K_{X/C}$。由 Hurwitz 公式 $K_{C'}+R'=g^*(K_C+R)$，得到 $K_{X'}+f'^{-1}(R')=h^*(K_X+f^{-1}(R))$，从而所证之主张成立。

由假设 $(X,f^{-1}(R)+B)$ 是 lc 的，故由上述主张与命题 5.20，$(X',f^{-1}(R')+B')$ 是 lc 的。对任意给定的 $c'\in C'$，可以选取 $R$ 使 $c'\in R'$，因此 $f'\colon(X',B')\to C'$ 是 lc 的。
\end{proof}

下面两个定理构成半稳定极小模型纲领的核心部分。终止定理在本节末尾证明，存在定理在 §7.4 证明。

\begin{theorem*}{7.7}
设 $(X,B)$ 是 Q-分解的三维 dlt 对，$B$ 有效，$f\colon X\to Y$ 是射影态射。又设 $g\colon Y\to C$ 是到非奇异曲线的
% [原文 p.261]
平坦态射，且 $g\circ f\colon(X,B)\to C$ 是 lc 的。此时，$Y$ 上任意一串射线型 $(K_X+B)$-翻转都是有限的。
\end{theorem*}

\begin{theorem*}{7.8}
（三维半稳定翻转的存在（参见推论 7.42））设 $(X,B)$ 是三维 Q-分解 dlt 对，$B$ 有效。设 $f\colon X\to Y$ 是射线型 $(K_X+B)$-翻转收缩。假设存在到非奇异曲线的平坦态射 $g\colon Y\to C$，使 $g\circ f$ 为 lc。此时，$f$ 有翻转。
\end{theorem*}

基于这些，便得到本章的主要定理。

\begin{theorem*}{7.9}
（半稳定极小模型定理）设给定三维 Q-分解对 $(X,B)$ 与非奇异曲线 $C$。设 $f\colon X\to Y$ 是射影态射，$g\colon Y\to C$ 是使 $g\circ f\colon(X,B)\to C$ 为 dlt 的平坦态射。

此时，$(X,B)$ 关于 $Y$ 的 MMP 的所有阶段都存在，且纲领终止。由此得到的对 $(X^m,B^m)$ 与态射 $f^m\colon X^m\to Y$ 具有以下性质。
\begin{enumerate}[label=(\arabic*),leftmargin=2.4em,itemsep=1pt,topsep=2pt,parsep=0pt]
\item 射 $g\circ f^m\colon(X^m,B^m)\to C$ 是 dlt 的。
\item 下列只有一个成立。
\begin{enumerate}[label=(\alph*),leftmargin=2.2em,itemsep=1pt,topsep=1pt,parsep=0pt]
\item $K_{X^m}+B^m$ 是 $f^m$-nef 的。
\item 存在 Fano 纤维空间的结构 $X^m\to Z^m\to Y$。
\end{enumerate}
\end{enumerate}
\end{theorem*}

% [Tutor Note: 7.9 比通常的 MMP 定理多出的是 (1)：MMP 全程保持 dlt，故可在相对设定下反复使用 7.7 与 7.42，而不必担心中间步骤退化出定义 7.1 未覆盖的纤维。]

上述定理在只假设 $(X,B)$ 是 dlt 且 $g\circ f$ 是 lc 时也成立。不过，此时 $g\circ f^m\colon(X^m,B^m)\to C$ 只能是 lc 的。

\begin{proof}
§3.7 说明了如何运行 MMP。由推论 3.44，dlt 条件得以保持。翻转由推论 7.42 存在，纲领由定理 7.7 终止。这就得到了 (1)。(2) 不过是 3.31 阶段 3 的特殊情形。
\end{proof}

\begin{theoremp}{7.10}
（半稳定标准模型定理）记号与假设同定理 7.9。再假设 $K_X+B$ 是 $f$-大的。此时，$(X,B)$ 在 $Y$ 上有标准模型。把它记为对 $(X^c,B^c)$ 与射 $f^c\colon X^c\to Y$，则 $g\circ f^c\colon(X^c,B^c)\to C$ 是 lc 的。
\end{theoremp}

\begin{proof}
只在 $(X,B)$ 是 klt 的情形给出证明。一般情形在最后讨论。

% [原文 p.262]
设 $(X^m,B^m)$ 与 $f^m\colon X^m\to Y$ 是极小模型（定理 7.9）。由推论 3.42 与推论 3.43，$(X^m,B^m)$ 是 klt 的。

由于 $K_{X^m}+B^m$ 是 $f^m$-nef 且 $f^m$-大的，由相对无基点定理 3.24，对 $r\gg0$，$r(K_{X^m}+B^m)$ 是 $f^m$-自由的。由此确定的态射的 Stein 分解记为
\[
f^m\colon X^m\xrightarrow{\ h\ }X^c\xrightarrow{\ f^c\ }Y,
\]
并记 $B^c:=h_*B^m$。于是
\[
K_{X^m}+B^m\equiv h^*(K_{X^c}+B^c),
\]
故由引理 2.30 可知 $(X^c,B^c)$ 是 klt 的，且 $g\circ f^c\colon(X^c,B^c)\to C$ 是 lc 的。

若只假设 $(X,B)$ 是 dlt，则一般的无基点定理不适用，必须注意 $\lfloor B^m\rfloor$ 来处理。这在 [K+92, 8.4] 中实施。其证明基于关于曲面 $\lfloor B^m\rfloor$ 的详细结果。因此只在三维得到证明。
\end{proof}

以下，本节致力于证明半稳定翻转的终止（定理 7.7）。下面叙述的定理 7.11 是其核心的局部版本。由引理 7.12，（定理 7.11）$\Rightarrow$（定理 7.7）是显然的。叙述这两个结果之后，再证明它们。

\begin{theorem*}{7.11}
设 $(X,B)$ 是三维 Q-分解 dlt 对，$f\colon X\to Y$ 是射影态射，而 $g\colon Y\to C$ 是到非奇异曲线的平坦态射。设 $T\subset Y$ 是 $g$-平凡的 Cartier 除子，且 $S=f^*T$。假设 $(X,S+B)$ 是 lc 的。此时，$Y$ 上任一串在 $Y\setminus T$ 上是同构的 $(K_X+B)$-翻转都是有限的。
\end{theorem*}

\begin{lemmap}{7.12}
记号同定理 7.7，但不假设 $X$ 的 Q-分解性。此时，存在 Zariski 开集 $C^\circ\subset C$，使得 $Y$ 上任一 $(K_X+B)$-翻转收缩 $h\colon X\to X^*$ 在 $(g\circ f)^{-1}(C^\circ)$ 上都是同构。
\end{lemmap}

\begin{proof}
这里给出两组证明。情形 1 给出直截了当、适用于许多情形的证明。情形 2 使用形变理论，概述一个适用于所有情形的证明。由于不使用 Q-分解性，这些证明稍加修改在解析情形也有效。

% [原文 p.263]
\textbf{情形 1}\quad 假设对适当的 $m>0$，$f_*O_X(m(K_X+B))\ne0$。这里先指出：若 $f$ 是双有理的，或 $K_X+B$ 在 $f$ 的一般纤维上是大的，则该假设成立。

为证明引理 7.12，只需对每个仿射开集 $U\subset Y$ 在 $f^{-1}(U)$ 上讨论即可（这是因为 $Y$ 被有限个这样的 $U$ 覆盖）。因此可以假设 $Y$ 是仿射的。于是存在 $X$ 上的有效 Cartier 除子 $V$，满足 $V\sim_f m(K_X+B)$。设 $V_i$ 是 $V$ 的不可约分支，$\bar V_i$ 是 $V_i$ 的正规化，$\bar V_i\to W_i$ 是 $\bar V_i\to Y$ 的 Stein 分解。

任取 Zariski 开集 $C^\circ\subset C$，使对所有 $i$，$\bar V_i|_{C^\circ}\to W_i|_{C^\circ}$ 都光滑且纤维维数不超过 1（由于 $\bar V_i$ 是正规曲面，这样的 $C^\circ$ 存在）。下面证明 $C^\circ\subset C$ 满足引理 7.12 的条件。

假设约化曲线 $\ell\subset\operatorname{Ex}(h)$ 满足 $P:=(g\circ f)(\ell)\in C^\circ$。由
\[
(V\cdot\ell)=m(K_X+B\cdot\ell)<0,
\]
对适当的 $a$ 有 $\ell\subset V_a$。这里 $f(\ell)$ 是一点，故 $\bar V_a|_{C^\circ}\to W_a|_{C^\circ}$ 有 1 维纤维，且作为集合，$\ell$ 与 $\bar V_a|_{C^\circ}\to W_a|_{C^\circ}$ 的某一个纤维的像完全一致。

由刚性引理 1.6，$\bar V_a|_{C^\circ}\to W_a|_{C^\circ}$ 的纤维在 $\bar V_a\to X\to X^*$ 下的像都是点。于是 $h$ 收缩 $V_a$，故它不是翻转收缩。

% [Tutor Note: 情形 1 的关键是刚性引理 1.6：$\bar V_a$ 的整个纤维在 $\bar V_a\to X\to X^*$ 下被压成点，故 $h$ 收缩掉 $V_a$，不可能是翻转收缩；情形 2 改用 $(-1)$-曲线的形变替代这一手。]

\textbf{情形 2}（一般情形的概略）\quad 设 $C^\circ$ 是 Zariski 开集，$X^\circ:=(g\circ f)^{-1}(C^\circ)$ 有奇点消解 $\pi\colon Z\to X^\circ$，使 $K_Z$ 是 $h$-nef 且 $Z\to C^\circ$ 光滑。此时证明 $C^\circ$ 满足引理 7.12 的条件。

设约化曲线 $\ell\subset\operatorname{Ex}(h)$ 满足 $P:=(g\circ f)(\ell)\in C^\circ$。导出矛盾。如情形 1 所论，可假设 $(\ell\cdot B)\ge0$。

存在适当的有效 $\pi$-例外 $\Q$-除子 $E$，使 $\pi^*K_{X^\circ}=K_Z+E$（引理 3.39）。对非奇异曲面 $Z_P$ 上的 $\ell':=\pi^{-1}_*(\ell)$，成立
\[
(\ell'\cdot K_{Z_P})=(\ell'\cdot K_Z)\le(\ell'\cdot\pi^*K_{X^\circ})=(\ell\cdot K_X)\le(\ell\cdot K_X+B)<0.
\]

% [原文 p.264]
于是可知 $\ell'$ 是 $Z_P$ 上的 $(-1)$-曲线，这是因为 $h|_{Z_P}\colon Z_P\to h(Z_P)$ 收缩 $\ell$。已知 $(-1)$-曲线 $\ell'$ 随 $P\in C^\circ$ 的形变而形变。（例如，把定理 1.10 证明的阶段 2 的方法应用于 $\ell'\to Z$，可知 $\ell'\to Z$ 的形变空间维数至少为 4，于是 $\ell'$ 运动，从而 $\ell\subset X$ 也运动。）由此，与情形 1 同样得到矛盾。
\end{proof}

\par\medskip\noindent\textbf{7.13}（定理 7.11 的证明）\quad 标准翻转的终止（定理 6.17）的证明稍加修改，便可作为定理 7.11 的证明。因此只给出所需的引理与定义。

\begin{lemmap}{7.14}
设 $X$ 是正规 $n$ 维簇，$S$ 是 Cartier 除子，且 $\Delta=\sum a_iD_i$，其中 $D_i$ 是互不相同的素除子，假设 $a_i\ge0$。假设 $(X,S+\Delta)$ 是 lc 的。如下定义：
\[
\Sigma_S(c)=\{\,X\ \text{上方的例外除子}\ E\mid \operatorname{center}_X(E)\subset S,\ a(E,X,\Delta)<c\,\}.
\]
此时，$\Sigma_S(0)$ 是空集。若 $(X,\Delta)$ 是 klt 的，则 $\Sigma_S(1)$ 是有限集。
\end{lemmap}

\begin{proof}
由 $a(E,X,\Delta)\ge a(E,X,S+\Delta)+1\ge0$ 得到第一个主张。

若 $(X,\Delta)$ 是 klt 的，则由推论 2.35，$(X,(1-\eps)S+\Delta)$ 也是 klt 的。因此，满足 $a(E,X,(1-\eps)S+\Delta)<0$ 的 $X$ 上方的例外除子 $E$ 只有有限个（命题 2.36）。若 $\operatorname{center}_X(E)\subset S$ 且 $a(E,X,(1-\eps)S+\Delta)\ge0$，则
\[
a(E,X,\Delta)\ge a(E,X,(1-\eps)S+\Delta)+(1-\eps)\ge1-\eps
\]
成立。若 $m(K_X+\Delta)$ 是 Cartier 的，则 $a(E,X,\Delta)$ 是有理数，其分母是 $m$ 的约数。于是，若 $0<\eps<1/m$，则 $a(E,X,\Delta)\ge1-\eps$ 与 $a(E,X,\Delta)\ge1$ 等价。
\end{proof}

\begin{definition}{7.15}
（参看定义 6.20）设 $(X,S+\Delta)$ 是 lc 对，$(X,\Delta)$ 是 dlt 对。其中 $S$ 是 Cartier 除子，$\Delta=\sum a_iD_i$ 是到互不相同的素除子 $D_i$ 的分解。记 $a=\max\{a_i\}$，$M=\sum a_i\Z_{\ge0}\subset\Q$。此时，定义
\[
d_S(X,\Delta)=\sum_{\xi\in M,\ \xi\ge a}\#\{\,E\mid \operatorname{center}_X(E)\subset S,\ a(E,X,\Delta)<1-\xi\,\},
\]
其中 $E$ 取遍 $X$ 上方的例外除子。

% [原文 p.265]
由引理 7.14，$d_S(X,\Delta)<\infty$。（注意：当 $a=1$ 时 $\Sigma_S(0)=\emptyset$；当 $a<1$ 时 $|\Sigma_S(1)|<\infty$。）此外，如定理 7.11 所示，$d_S(X,\Delta)$ 在 $(K_X+S+\Delta)$-翻转下不增加（引理 3.38）。
\end{definition}

\begin{lemma}{7.16}
（参看引理 6.21）设 $\phi\colon X\dashrightarrow X'$ 是定理 7.11 中所述的 $(K_X+S+\Delta)$-翻转，其中 $\Delta=\sum a_iD_i$。设 $C'\subset X'$ 是后翻转曲线，$E_{C'}$ 是在 $C'$ 的泛点附近爆破 $C'$ 所得的例外除子。此时，$X'$ 沿 $C'$ 非奇异，且成立
\[
0\le a(E_{C'},X,\Delta)<a(E_{C'},X',\phi_*\Delta)=1-\sum a_i\cdot\operatorname{mult}_{C'}(D_i).
\]
\end{lemma}

这些修正后的引理与定义一经采用，定理 6.17 的证明便成为定理 7.11 的证明。\hfill$\blacksquare$

\section{半稳定约化定理}
\takeaway{本节在三维情形证明半稳定约化定理：先取分歧覆盖使纤维约化，再用 Q-分解终端化消解新生的环面奇点，从而得到 snc 模型。}

本节相当于推论 7.42 证明的第一阶段。这一阶段称为\textbf{半稳定约化定理}（semi-stable reduction theorem），在所有维数都已证明。

$X$ 是 $\C$ 上的正规代数簇，或正规解析空间的紧致部分集合的适当邻域。$f\colon X\to C$ 是到 $\C$ 上非奇异曲线 $C$ 的平坦态射，$B\subset X$ 是闭子集。对任意态射 $C'\to C$，记 $X':=X\times_CC'$，以 $X'_n$ 表示 $X'$ 的正规化。对态射 $X''\to X'_n$，用 $B''\subset X''$ 表示由诱导态射 $X''\to X$ 拉回的 $B$。

\begin{theorem*}{7.17}
（半稳定约化定理 [KKMSD73]）记号如上。存在由非奇异曲线 $C'$ 出发的有限态射 $C'\to C$ 与射影奇点消解 $g\colon X''\to X'_n$ 组成的组，使诱导态射 $f''\colon X''\to C'$ 具有以下性质。
\begin{enumerate}[label=(\arabic*),leftmargin=2.4em,itemsep=1pt,topsep=2pt,parsep=0pt]
\item 对每个 $c'\in C'$，$(f'')^*(c')\cup\operatorname{Ex}(g)\cup B''$ 是 snc 除子。
\item 对 $c'\in C'$，$(f'')^*(c')$ 是约化的。
\end{enumerate}
\end{theorem*}

这比 [KKMSD73] 中的定理稍为精确。下面说明其证明如何进行。

由定理 0.2 与定理 0.3，可进行爆破（的合成）$g\colon X_1\to X$，使 $X_1$ 非奇异，且
\[
\operatorname{Ex}(g)\cup g^{-1}(B)\cup f_1^{-1}(c)
\]
对所有 $c\in C$ 都是 snc 除子，这样便得到 $f_1\colon X_1\to C$。

% [原文 p.266]
然而在大多数情形下，$f_1^{-1}(c)$ 不是约化的。尽管如此，若对 $f_1\colon X_1\to C$ 与闭子集 $B_1:=\operatorname{Ex}(g)\cup g^{-1}(B)$ 证明了定理 7.17，则对 $f\colon X\to C$ 与 $B$ 定理 7.17 显然成立。因此，可以把 $X$ 换成 $X_1$，把 $B$ 换成 $B_1$。特别地，可以假设 $X$ 非奇异，且 $f^*(c)\cup B$ 对所有 $c\in C$ 是 snc 的。

下一阶段是取分歧覆盖 $C'\to C$。由此容易使 $X'_n\to C'$ 的所有纤维都成为约化的。$X'_n$ 会出现奇点，但这些奇点全是环面奇点（参看命题 7.23）。关键在于非常小心地消解这些奇点。[KKMSD73, Theorem 11*] 包含 $B$ 在内证明了一切。但那需要环面嵌入的组合理论。

本节的目标是：至少对三维情形，证明 $X'_n$ 的 Q-分解终端化给出定理所主张的\textbf{半稳定奇点消解}（semi-stable resolution），而不使用环面嵌入理论。

\begin{remark}{7.18}
若使用所需形式的终端化（定理 6.23）与 Q-分解化（定理 6.25），则下述定理 7.19 在代数设定与解析设定下同样有效。
\end{remark}

\begin{theorem*}{7.19}
（三维簇的半稳定约化定理）记号同定理 7.17。设 $X$ 是非奇异三维簇，$B\subset X$ 是除子，且假设对所有 $c\in C$，$f^{-1}(c)+B$ 是 snc 除子。设 $p\colon C'\to C$ 是从非奇异曲线出发的有限态射，满足以下性质。

对任意 $c\in C$、$c'\in p^{-1}(c)$，$f^*(c)$ 的任意不可约分支的重数都是 $p$ 在 $c'$ 处分歧指数的约数。

此时，正规化 $X'_n$ 只有标准商奇点。此外，取 $X'_n$ 的非奇异射影终端化 $X'_{tf}$（参看定理 6.23、定理 6.25），则 $X'':=X'_{tf}$ 满足定理 7.17 的结论。
\end{theorem*}

% [Tutor Note: 半稳定约化的含义是：先允许有限分歧覆盖把纤维重数化成 1，再用射影终端化把边界整理成 snc。这样族中纤维的退化被化为组合数据，因而在任意维数都能陈述。]

\begin{remark}{7.20}
定理 7.19 及其证明是二维版本 [DM69] 的自然推广。不过在二维版本中，「非奇异射影终端化」就是极小奇点消解。
\end{remark}
